% GENERATED FILE. Edit canonical lessons, then run npm run generate:books.
% canonical-source-hash: fnv1a64:17d16bcbf5ade7e8
% canonical-lessons: ES-C17-sintesis-futuro-condicional

\chapter{Synthesis --- Promising, and Supposing}
\label{ch:spanish-synthesis-future-conditional}
\hlchaptermodality{120}

\begin{chapteropening}
\textbf{By the end of this chapter:} \emph{I can choose between a promise and a supposition, and say why the endings differ.}

Haber glued on twice - present for the promise, past for the guess.
\end{chapteropening}

\section[Practice]{Synthesis — promising, and supposing}

\label{lesson:ES-C17-sintesis-futuro-condicional}



\begin{quote}
Nothing new in this chapter. Say \emph{hablaré}, then \emph{hablaría}.
\end{quote}

\begin{grammarlens}[title={Grammar Lens: the two endings are one word, twice}]
Three chapters ago you were told where these endings came from. Now it is worth seeing what that fact \emph{buys} you.

Both are the verb \textbf{haber} — Latin \emph{habēre}, \textquotedblleft{}to have\textquotedblright{} — glued onto the back of the infinitive. Romance speakers stopped saying \textquotedblleft{}I have to speak\textquotedblright{} as two words and fused it into one:

\begin{quote}
\emph{hablar} + \textbf{he} (I have) $\to$ \textbf{hablaré} — I will speak. \emph{hablar} + \textbf{había} (I had) $\to$ \textbf{hablaría} — I would speak.
\end{quote}

That is the entire difference between the two tenses: \textbf{the same auxiliary, once in the present and once in the past.} Not two systems to learn — one system, looked at from two moments.
\end{grammarlens}

\begin{grammarlens}[title={What you've built}]
And the meaning follows from the grammar rather than being extra to it.

\emph{Having to} do something \textbf{now} is a commitment: \textbf{hablaré} is a promise, or a prediction you are standing behind. \emph{Having had to} do it is a supposition — the obligation sits in a past that never arrived, so \textbf{hablaría} is hypothetical.

\begin{quote}
Promise from the present. Suppose from the past.
\end{quote}

English hides the same move: \emph{will} is the present of an old verb meaning \emph{to want}, and \emph{would} is its past. \textbf{Will/would} is the same pair as \textbf{-é/-ía}.
\end{grammarlens}

\subsection*{The exchange}
\begin{quote}
— ¿\textbf{Hablarás} español en Madrid? — Sí, \textbf{hablaré} español.
\end{quote}

>

\begin{quote}
— ¿Y sin práctica? — \textbf{Hablaría} español, pero no bien.
\end{quote}

One letter-run apart, and the second speaker has moved from a commitment to a guess without changing a single other word.

\subsection*{Your turn}
\begin{itemize}
\raggedright
  \item \emph{Say it:} \emph{hablaré} then \emph{hablaría}, and hear a promise turn into a guess
\end{itemize}

\emph{Twice through:} Until the choice is about what you mean, not which ending it was.

\begin{tcolorbox}[breakable,colback=teal!4,colframe=teal!35!black,title={Before you move on}]
What is glued onto the infinitive in both tenses? (\textbf{Haber} — the verb \emph{to have}.) Which one is a promise, and which a supposition? (\emph{-é} the promise, \emph{-ía} the supposition.) And what is the English pair that works the same way? (\textbf{Will} and \textbf{would} — a present and a past of one old verb.)
\end{tcolorbox}
